\documentclass{article}

\usepackage[preprint]{neurips_2025}

\usepackage[utf8]{inputenc}
\usepackage[T1]{fontenc}
\usepackage{hyperref}
\usepackage{url}
\usepackage{booktabs}
\usepackage{amsfonts}
\usepackage{nicefrac}
\usepackage{microtype}
\usepackage{xcolor}
\usepackage{graphicx}
\usepackage{listings}
\usepackage{subcaption}

\lstset{
  basicstyle=\ttfamily\small,
  breaklines=true,
  frame=single,
  language=Go,
  keywordstyle=\color{blue},
  commentstyle=\color{gray},
  stringstyle=\color{red},
}

\title{SWE-Bench 5G: A Benchmark for Evaluating AI Coding Agents \\on 5G Core Network Engineering Tasks}

\author{
  Anonymous Author(s)
}

\begin{document}

\maketitle

\begin{abstract}
Existing benchmarks for AI coding agents---such as SWE-Bench, SWE-Bench Mobile, and BeyondSWE---have primarily targeted general-purpose software in Python, Swift, or cross-repository settings. However, none address the \emph{telecommunications domain}, where software correctness is governed by formal protocol specifications (3GPP), systems are inherently distributed, and bugs often manifest through complex protocol state machine interactions. We introduce \textbf{SWE-Bench 5G}, a benchmark that evaluates coding agents on real bug-fixing tasks drawn from \textbf{free5GC}, the most widely adopted open-source 5G core network implementation. Each task instance is packaged as a self-contained Docker image with a pre-installed Go toolchain, the source code at the buggy commit, existing regression tests, and fail-to-pass tests that verify the fix. Our benchmark introduces three challenges absent from prior work: (1) \emph{specification-driven reasoning}, where agents must interpret 3GPP technical specifications to understand expected behavior; (2) \emph{distributed network function coordination}, where bugs span multiple interacting microservices; and (3) \emph{protocol state machine complexity}, where fixes require understanding stateful signaling procedures. We release 10 validated task instances spanning 7 network functions, with an automated pipeline for scaling to 50+ instances from a candidate pool of 280 issues mined across 16 repositories.
\end{abstract}

\section{Introduction}

The rapid advancement of large language models (LLMs) has given rise to a new generation of \emph{AI coding agents}---systems that autonomously navigate codebases, diagnose bugs, and produce patches \citep{jimenez2024swebench}. Benchmarks such as SWE-Bench \citep{jimenez2024swebench}, SWE-Bench Mobile \citep{yao2025swebenchmobile}, and BeyondSWE \citep{shi2025beyondswe} have been instrumental in measuring progress. Yet these benchmarks share a critical blind spot: they evaluate agents exclusively on general-purpose software, leaving open the question of whether coding agents can operate effectively in \emph{vertical, specification-driven domains}.

The 5G telecommunications sector presents a compelling test case. The 5G core network (5GC) is specified by the 3rd Generation Partnership Project (3GPP) across hundreds of Technical Specifications (TS), each defining protocol messages, state machines, and service-based interfaces with precise semantics. Implementations such as free5GC\footnote{\url{https://github.com/free5gc/free5gc}} translate these specifications into production-grade Go code, decomposed into multiple Network Functions (NFs)---AMF, SMF, PCF, UPF, and others---that communicate via HTTP/2 and PFCP protocols. Bugs in this domain are qualitatively different from those in typical open-source projects:

\begin{itemize}
    \item \textbf{Specification-driven correctness.} A ``correct'' fix is not merely one that passes tests, but one that conforms to the relevant 3GPP TS clauses. Agents must locate, interpret, and apply normative protocol text.
    \item \textbf{Distributed microservice architecture.} Network Functions run as independent processes. A single user operation (e.g., UE registration) triggers a signaling cascade across AMF, AUSF, UDM, SMF, and UPF. Bugs may require understanding cross-NF interactions.
    \item \textbf{Protocol state machine complexity.} 5G procedures (registration, session establishment, handover) are defined as multi-step state machines with specific message sequences, timer constraints, and error handling paths.
\end{itemize}

We present \textbf{SWE-Bench 5G}, the first benchmark targeting AI coding agents in the telecommunications domain. Following the BeyondSWE methodology of providing fully reproducible Docker environments, each task instance packages a complete Go development environment with the free5GC source code checked out at the buggy commit, along with fail-to-pass tests that validate the fix. Our contributions are:

\begin{enumerate}
    \item A benchmark framework that maps 5G core network bugs to the SWE-Bench task format, enriched with 3GPP specification references and protocol context.
    \item A dual test strategy---\emph{direct function call} and \emph{diff-based intent testing}---tailored to the distributed, specification-driven nature of 5G software where many functions cannot be unit-tested in isolation.
    \item A dataset of 10 validated task instances spanning 7 NFs, drawn from real free5GC issues, with an automated pipeline for expansion across 16 repositories (280 candidates identified).
    \item An end-to-end evaluation harness and batch tooling for building, validating, and scoring agent performance on task instances.
\end{enumerate}

\section{Related Work}

\subsection{Software Engineering Benchmarks for Coding Agents}

\textbf{SWE-Bench} \citep{jimenez2024swebench} pioneered the evaluation of coding agents on real-world GitHub issues, drawing 2,294 task instances from 12 popular Python repositories. Each instance pairs an issue description with a fail-to-pass test, enabling automated evaluation. Subsequent work has extended this paradigm along several axes.

\textbf{SWE-Bench Mobile} \citep{yao2025swebenchmobile} targets iOS development with multi-modal inputs (PRDs and Figma designs) over a 500K-line Swift/Objective-C codebase. It introduces \emph{diff-based intent tests} that verify patch correctness through structural inspection rather than runtime execution---an approach we adapt for 5G functions with complex runtime dependencies.

\textbf{BeyondSWE} \citep{shi2025beyondswe} expands the resolution scope to cross-repository reasoning, domain-specific problem solving, and dependency-driven migrations, packaging each of its 500 instances in a Docker image for full reproducibility. We adopt the same Docker-based environment design.

Despite this diversity, all existing benchmarks target general-purpose software. None evaluate agents on \emph{protocol-specified, distributed telecommunications systems}---a domain with distinct challenges around specification compliance, signaling procedures, and network function coordination.

\subsection{Open-Source 5G Core Networks}

free5GC \citep{free5gc} is an open-source implementation of the 3GPP Release 15+ 5G core network, written in Go. It implements the full set of control-plane NFs (AMF, SMF, PCF, UDM, UDR, AUSF, NRF, NSSF, N3IWF) and a user-plane function (UPF). With 2.3K GitHub stars and active development, it is the most widely used open-source 5GC for research and prototyping. Its Go codebase and modular architecture (one repository per NF) make it amenable to automated code analysis and per-NF testing.

\section{Benchmark Design}

\subsection{Task Formulation}

Following the SWE-Bench paradigm, we define each task instance as a tuple $\mathcal{T} = (\mathcal{I}, \mathcal{O}, \mathcal{E})$:

\begin{itemize}
    \item \textbf{Input} $\mathcal{I}$: A problem statement derived from the original GitHub issue, augmented with references to relevant 3GPP Technical Specification clauses. The agent also has access to the full source code of the affected NF(s) at the buggy commit.
    \item \textbf{Output} $\mathcal{O}$: A unified diff patch that resolves the issue.
    \item \textbf{Evaluation} $\mathcal{E}$: A test suite partitioned into \texttt{PASS\_TO\_PASS} tests (existing tests that must continue passing) and \texttt{FAIL\_TO\_PASS} tests (new tests that fail on the buggy version and must pass after the fix).
\end{itemize}

\subsection{Task Taxonomy}

We categorize tasks along two dimensions: \emph{scope} and \emph{bug type}.

\paragraph{Scope.}
\begin{itemize}
    \item \textbf{SingleNF}: The bug and fix are contained within a single Network Function (all 10 current instances).
    \item \textbf{CrossNF}: The fix requires coordinated changes across multiple NFs (planned).
    \item \textbf{Protocol}: The bug violates a specific 3GPP protocol requirement (planned).
    \item \textbf{DataPlane}: The bug affects user-plane forwarding in the UPF or PFCP signaling (planned).
\end{itemize}

\paragraph{Bug type.} Our current dataset covers: \emph{nil pointer dereference} (6 instances), \emph{crash / index out of range} (2), and \emph{missing validation} (2).

\subsection{Dual Test Strategy}

A key methodological contribution is our \emph{dual test strategy}, developed to address the challenge that many 5G NF functions cannot be called directly in unit tests due to complex runtime dependencies (SCTP connections, MongoDB contexts, inter-NF signaling).

\paragraph{Strategy A: Direct Function Call.} When the buggy function can be invoked with minimal setup, we call it directly with crash-triggering inputs and use Go's \texttt{defer/recover} mechanism to detect panics:

\begin{lstlisting}[caption={Strategy A: Direct function call (pilot\_pcf\_879)}]
func TestNilRouteReq(t *testing.T) {
    defer func() {
        if r := recover(); r != nil {
            t.Fatalf("BUG: panic: %v", r)
        }
    }()
    // Call the actual buggy function
    provisioningOfTrafficRoutingInfo(smPolicy, "app1", nil, "")
}
\end{lstlisting}

\paragraph{Strategy B: Diff-Based Intent Test.} Inspired by SWE-Bench Mobile's structural patch inspection \citep{yao2025swebenchmobile}, when the function requires complex runtime context, we verify that the \emph{source code} contains the necessary fix pattern:

\begin{lstlisting}[caption={Strategy B: Diff-based intent test (amf\_pr161)}]
func TestHandlerHasGNbIdNilCheck(t *testing.T) {
    data, _ := os.ReadFile("handler.go")
    if !strings.Contains(string(data), "GNbId != nil") {
        t.Fatal("BUG: accesses GNbId without nil check")
    }
}
\end{lstlisting}

Strategy B ensures that agents must modify the actual source code to pass the test, while avoiding the need to construct complex mock objects for NGAP connections, SCTP sessions, or database contexts.

\begin{table}[h]
\centering
\caption{Test strategy selection criteria and usage}
\label{tab:strategies}
\begin{tabular}{lcp{5cm}c}
\toprule
\textbf{Strategy} & \textbf{Count} & \textbf{When to use} & \textbf{Example} \\
\midrule
A: Direct Call  & 3 & Function callable with simple args & pilot\_pcf\_879 \\
B: Diff-Based   & 7 & Complex dependencies (NGAP, DB) & amf\_pr161 \\
\bottomrule
\end{tabular}
\end{table}

\subsection{Docker Environment Design}

Following BeyondSWE \citep{shi2025beyondswe}, each task instance is packaged as a self-contained Docker image built from a shared \texttt{golang:1.25-bookworm} base. The image contains:

\begin{lstlisting}[language=bash, caption={Docker image internal structure}]
/opt/free5gc-<nf>/             # NF source at buggy commit
  internal/sbi/processor/      # Bug location
  go.mod, go.sum               # Dependencies (pre-downloaded)
/opt/test-suite/
  existing_test.go             # PASS_TO_PASS tests
  fail_test.go                 # FAIL_TO_PASS tests
  run_tests.sh                 # Test runner script
/opt/task/
  problem_statement.md         # Issue + 3GPP spec references
\end{lstlisting}

\subsection{Dataset Construction Pipeline}

Our automated pipeline consists of five stages:

\paragraph{Stage 1: Candidate mining.} An automated script (\texttt{mine\_issues.py}) scans 16 free5GC sub-repositories via the GitHub API, collecting closed issues with linked merged PRs. It filters for bug-fix PRs by title keywords and excludes refactoring, documentation, and dependency changes. This yielded \textbf{280 candidates}.

\paragraph{Stage 2: Quality filtering.} We retain candidates with clear issue descriptions, linked PRs that modify source code files, and bugs reproducible without external hardware. From 280 candidates, 29 met all quality criteria.

\paragraph{Stage 3: Commit identification.} For each candidate, we identify the \emph{parent commit} (buggy version) and \emph{fix commit} using GitHub API data, verified against the actual PR diff.

\paragraph{Stage 4: Test suite construction.} We construct fail-to-pass tests using Strategy A or B depending on function testability. Tests are authored based on the issue description, PR diff, and 3GPP specification requirements. A generation script (\texttt{generate\_instances.py}) templates the process.

\paragraph{Stage 5: Automated validation.} A batch script (\texttt{batch\_build\_validate.sh}) builds Docker images and runs three-step validation for each instance:
\begin{enumerate}
    \item At the parent commit: existing tests PASS, fail-to-pass tests FAIL.
    \item Apply the ground-truth fix (via \texttt{git cherry-pick} or \texttt{checkout}): all tests PASS.
\end{enumerate}
Only instances passing both steps are included. The current v0.2 release contains \textbf{10 validated instances} from 21 constructed.

\section{Dataset Statistics}

\begin{table}[h]
\centering
\caption{SWE-Bench 5G v0.2 dataset statistics}
\label{tab:dataset_stats}
\begin{tabular}{ll}
\toprule
\textbf{Attribute} & \textbf{Value} \\
\midrule
Total validated instances  & 10 \\
Source project             & free5GC (Go, Apache-2.0) \\
NFs represented            & 7 (AMF, PCF, SMF, UDM, NRF, NSSF, AUSF) \\
Bug types                  & nil pointer (6), crash (2), missing validation (2) \\
Difficulty distribution    & 10 Easy \\
Test strategy              & 3 Direct Call (A) + 7 Diff-Based (B) \\
Avg. tests per instance    & 3.2 (1.2 existing + 2.0 fail-to-pass) \\
3GPP specs referenced      & TS 24.501, 29.507, 29.509, 29.510, 29.512, \\
                           & 29.514, 29.518, 29.531, 38.413 \\
Candidate pool             & 280 (from 16 repositories) \\
Docker base image          & \texttt{golang:1.25-bookworm} ($\sim$1.2 GB shared) \\
\bottomrule
\end{tabular}
\end{table}

\begin{table}[h]
\centering
\caption{Validated task instances in v0.2}
\label{tab:instances}
\begin{tabular}{llllcc}
\toprule
\textbf{Instance} & \textbf{NF} & \textbf{Bug Type} & \textbf{Strategy} & \textbf{Tests} & \textbf{Spec} \\
\midrule
pilot\_pcf\_879 & PCF  & nil pointer     & A & 5 & TS 29.514 \\
amf\_pr118      & AMF  & crash           & B & 2 & TS 24.501 \\
amf\_pr157      & AMF  & missing valid.  & A & 3 & TS 24.501 \\
amf\_pr161      & AMF  & nil pointer     & B & 3 & TS 38.413 \\
ausf\_pr52      & AUSF & nil pointer     & B & 2 & TS 29.509 \\
nrf\_pr79       & NRF  & crash           & B & 2 & TS 29.510 \\
nssf\_pr39      & NSSF & nil pointer     & B & 3 & TS 29.531 \\
pcf\_pr57       & PCF  & nil pointer     & B & 2 & TS 29.507 \\
smf\_pr125      & SMF  & nil pointer     & B & 2 & TS 29.518 \\
udm\_pr45       & UDM  & index OOB       & B & 3 & TS 29.509 \\
\bottomrule
\end{tabular}
\end{table}

\subsection{NF Coverage}

\begin{table}[h]
\centering
\caption{Network Function coverage in current dataset}
\label{tab:nf_coverage}
\begin{tabular}{llc}
\toprule
\textbf{NF} & \textbf{Role} & \textbf{Instances} \\
\midrule
AMF  & Access \& Mobility Management & 3 \\
PCF  & Policy Control                & 2 \\
SMF  & Session Management            & 1 \\
UDM  & Unified Data Management       & 1 \\
NRF  & NF Repository / Discovery     & 1 \\
NSSF & Network Slice Selection       & 1 \\
AUSF & Authentication Server         & 1 \\
\midrule
UPF  & User Plane (planned)          & 0 \\
\bottomrule
\end{tabular}
\end{table}

\section{Case Studies}

\subsection{Strategy A: PCF PolicyAuthorization (pilot\_pcf\_879)}

\paragraph{Bug.} The PCF crashes with nil pointer dereference when \texttt{suppFeat="1"} but \texttt{AfRoutReq} is absent. The function \texttt{provisioningOfTrafficRoutingInfo()} unconditionally dereferences its \texttt{routeReq} parameter.

\paragraph{Test design.} We call the function directly with \texttt{routeReq=nil}. On the buggy version, Go's \texttt{recover()} catches the panic and marks the test as FAIL. After the fix adds a nil check, the function returns gracefully and the test passes.

\paragraph{3GPP context.} The \texttt{suppFeat} field (TS 29.514 bit 1) negotiates traffic-routing support. The specification does not mandate \texttt{AfRoutReq} presence---the PCF should handle its absence gracefully.

\subsection{Strategy B: AMF GNbId Nil Check (amf\_pr161)}

\paragraph{Bug.} The AMF crashes when \texttt{handleUplinkRANConfigurationTransferMain} accesses \texttt{targetRanNodeID.GNbId.GNBValue} without nil-checking \texttt{GNbId}.

\paragraph{Challenge.} The handler function requires an active SCTP connection, NGAP context, and \texttt{AmfRan} state---impractical to mock in a unit test.

\paragraph{Test design.} We use diff-based intent testing: the test reads \texttt{handler.go} and verifies the string \texttt{"GNbId != nil"} appears in the source. On the buggy version, this pattern is absent (FAIL). After the agent adds the nil guard, the pattern is found (PASS). This approach forces the agent to modify the actual source code while avoiding complex dependency setup.

\section{Baseline Experiments}

\subsection{Experimental Setup}

We evaluate on the 10 validated instances with multiple frontier coding agents.

\paragraph{Agent configurations.}
\begin{itemize}
    \item Claude Code (Claude Opus 4.6)
    \item Cursor Agent (Claude Sonnet 4.6)
    \item OpenAI Codex CLI (GPT-4.1)
\end{itemize}

\paragraph{Evaluation metric.} Following SWE-Bench, we use the binary \emph{resolve rate}: an instance is \emph{resolved} if and only if (1) all \texttt{FAIL\_TO\_PASS} tests pass and (2) all \texttt{PASS\_TO\_PASS} tests continue to pass after applying the agent's patch.

\subsection{Results}

\begin{table}[h]
\centering
\caption{Baseline results on 10 validated instances (to be populated).}
\label{tab:results}
\begin{tabular}{lccc}
\toprule
\textbf{Agent} & \textbf{Model} & \textbf{Resolve Rate} & \textbf{Avg. Time (s)} \\
\midrule
Claude Code   & Opus 4.6    & TBD & TBD \\
Cursor Agent  & Sonnet 4.6  & TBD & TBD \\
Codex CLI     & GPT-4.1     & TBD & TBD \\
\bottomrule
\end{tabular}
\end{table}

\textit{Note: Baseline results will be populated upon completion of evaluation runs. We anticipate that easy, single-NF tasks will have moderate resolve rates, while the domain-specific nature of the bugs (3GPP specification knowledge, Go idioms for nil handling) will present challenges distinct from general-purpose benchmarks.}

\section{Challenges and Lessons Learned}

Our experience constructing SWE-Bench 5G reveals several challenges that distinguish it from existing benchmarks:

\paragraph{Zero existing test coverage.} Unlike mature Python projects that SWE-Bench draws from, free5GC has virtually no unit tests---most NF repositories contain zero \texttt{\_test.go} files. Every fail-to-pass test must be authored from scratch, making test quality the primary bottleneck for dataset scaling.

\paragraph{Test strategy selection.} A naive approach of directly calling buggy functions fails for functions requiring complex runtime context (NGAP sessions, database connections). Our dual strategy (direct call vs. diff-based intent) resolves this but introduces a trade-off: diff-based tests are less rigorous than functional tests, as they verify code patterns rather than runtime behavior.

\paragraph{Commit ancestry precision.} Identifying the correct ``buggy'' parent commit is non-trivial. Some PRs are based on branches that already contain partial fixes, and merge commits may introduce unrelated changes. We encountered cases where the parent commit already contained the fix pattern, requiring manual verification and re-selection of commits.

\paragraph{Environment reproducibility.} Docker images depend on external Go modules (\texttt{go mod download}) and base images from Docker Hub. Network-restricted environments (common in China) require mirrors for both (\texttt{goproxy.cn}, Docker registry mirrors), adding deployment complexity.

\paragraph{Domain knowledge scarcity.} Telecommunications is underrepresented in LLM training data relative to web development or data science. We hypothesize that agents will struggle more with 5G-specific tasks not because of code complexity, but because of unfamiliarity with the domain's concepts and conventions.

\section{Discussion and Future Work}

\paragraph{Scaling to 50+ instances.} Our candidate pool of 280 issues, combined with the automated mining and batch validation pipeline, provides a clear path to scaling. The primary bottleneck is test authoring quality---we plan to use LLM-assisted test generation with human review.

\paragraph{Medium and hard difficulty tasks.} The current release contains only easy (single-file, single-NF) tasks. We have identified 2 medium-difficulty candidates requiring multi-file changes and plan to expand to hard tasks involving cross-NF coordination.

\paragraph{3GPP specification augmentation.} Following SWE-Bench Mobile's multi-modal approach with Figma designs, we plan to augment task instances with 3GPP signaling flow diagrams and specification excerpts, testing whether providing domain context improves resolve rates.

\paragraph{SearchSWE-5G.} Inspired by BeyondSWE's SearchSWE framework, we plan to develop a 3GPP specification search tool (via MCP) that agents can query during debugging, enabling evaluation of specification-grounded reasoning.

\paragraph{Agent baseline analysis.} Beyond resolve rates, we plan to analyze agent behavior: Do agents read the 3GPP references in the problem statement? Do they explore related NF code? How do different prompt strategies (e.g., SWE-Bench Mobile's ``defensive programming'' prompt) affect performance on 5G tasks?

\section{Conclusion}

We have presented SWE-Bench 5G, the first benchmark for evaluating AI coding agents on telecommunications software engineering tasks. By targeting the free5GC open-source 5G core network, our benchmark introduces challenges around specification-driven reasoning, distributed network function coordination, and protocol state machine complexity that are absent from existing benchmarks. Our dual test strategy---combining direct function calls with diff-based intent testing---addresses the practical challenge of testing functions with complex runtime dependencies. The v0.2 release contains 10 validated instances spanning 7 network functions, with an automated pipeline supporting expansion from a 280-candidate pool. The dataset, evaluation harness, and project showcase are publicly available.

\begin{itemize}
    \item Dataset: \url{https://huggingface.co/datasets/tenderzada/SWEBench5G}
    \item Code: \url{https://github.com/tenderzada/swebench5g}
    \item Project page: \url{https://tenderzada.github.io/swebench5g/}
\end{itemize}

\begin{ack}
We thank the free5GC open-source community for their work on building and maintaining a high-quality 5G core network implementation. This work builds upon their contributions under the Apache-2.0 license.
\end{ack}

\bibliographystyle{plainnat}
\bibliography{references}

\end{document}
